En este capítulo se sintetizan las principales aportaciones del trabajo y se evalúa el
grado de consecución de los objetivos planteados en el Capítulo~1. La
Sección~\ref{sec:conclusiones} recoge las conclusiones extraídas del desarrollo e
implementación del sistema. La Sección~\ref{sec:lineas_futuras} identifica las líneas
de trabajo futuro que podrían ampliar las capacidades de Voctomix~2.0 o adaptar la
plataforma a nuevos casos de uso.

\section{Conclusiones}
\label{sec:conclusiones}

Este trabajo ha desarrollado y validado Voctomix~2.0, un sistema modular de producción
audiovisual en directo de código abierto que extiende la herramienta Voctomix original
con funcionalidades de nivel profesional. A continuación se evalúa el grado de
consecución de cada uno de los seis objetivos específicos planteados.

\vspace{0.3cm}
\noindent\textbf{OE-1: Funcionalidades de producción profesional.}
Se han implementado las tres funcionalidades broadcast que el sistema requería. La
rotulación dinámica permite superponer gráficos PNG y texto en tiempo real, con
fundido de entrada y salida configurable y desactivación automática en cada corte de
cámara. El gestor de estados de emisión (\textit{stream blanker}) ofrece tres modos
(LIVE, PAUSE y NOSTREAM) que permiten interrumpir la señal de salida sin detener el
procesamiento interno de la mezcla. La funcionalidad \textit{Audio Follows Video}
ajusta automáticamente los volúmenes en cada cambio de fuente, eliminando la gestión
manual del audio y reduciendo la carga cognitiva del operador durante la producción.

\vspace{0.2cm}
\noindent\textbf{OE-2: Servicio de telemetría.}
El servicio de telemetría registra en tiempo real los eventos de operación del mezclador
(tipo CHANGE) y el estado periódico del sistema (tipo STATE, cada 5 segundos) en
formato JSON Lines. Los eventos se publican simultáneamente en un broker RabbitMQ
mediante AMQP, lo que permite su integración con sistemas externos de monitorización
o análisis sin modificar el núcleo de la aplicación. Los datos de rendimiento
obtenidos a través de este sistema han permitido caracterizar cuantitativamente el
comportamiento del sistema en los cuatro escenarios de despliegue.

\vspace{0.2cm}
\noindent\textbf{OE-3: Contenerización con Docker.}
El sistema completo se ha empaquetado en una pila de diez contenedores Docker orquestada
con Docker Compose, incluyendo voctocore, cuatro fuentes de cámara sintéticas, el stream
blanker, el gestor de audio, el broker RabbitMQ y el servicio de telemetría. El despliegue
completo se realiza con un único comando (\texttt{docker compose up}), eliminando la
principal barrera de adopción del sistema.

\vspace{0.2cm}
\noindent\textbf{OE-4: Despliegue en Kubernetes.}
El sistema se desplegó sobre un clúster local Minikube mediante manifiestos declarativos
que reproducen la misma funcionalidad que el escenario Docker Compose en una
infraestructura orquestada. El orden de arranque controlado y las sondas de preparación
garantizan que las dependencias entre servicios se respetan automáticamente. La
arquitectura del sistema escala de forma natural desde Docker Compose hacia Kubernetes
sin necesidad de modificar el código del núcleo ni de la interfaz gráfica.

\vspace{0.2cm}
\noindent\textbf{OE-5: Validación en múltiples escenarios.}
El sistema ha sido validado funcionalmente en cuatro escenarios de despliegue
(monopuesto nativo, dos PCs en red local, Docker Compose y Kubernetes), superando
en todos ellos las nueve categorías de pruebas funcionales definidas. La latencia
de conmutación medida en los escenarios contenerizados fue inferior a 5\,ms en
el 100\,\% de las mediciones, con una mediana de 1{,}5\,ms en Docker y 1{,}8\,ms
en Kubernetes, cumpliendo con amplio margen el requisito de conmutación
imperceptible a 25 fotogramas por segundo (40\,ms por fotograma).

\vspace{0.2cm}
\noindent\textbf{OE-6: Despliegue en producción real.}
El sistema fue desplegado en un entorno de producción real dentro del proyecto
europeo de investigación \textbf{CyberNEMO} (UPM), donde actuó como motor de
producción audiovisual. Este despliegue acredita la madurez del sistema más allá
del prototipo académico y lo posiciona como una contribución directamente aplicable
a infraestructuras de investigación y producción audiovisual profesional.

\vspace{0.4cm}
La tabla~\ref{tab:objetivos_cumplidos} resume el grado de consecución de los objetivos
y las secciones de la memoria donde se documentan.

\begin{table}[H]
  \centering
  \caption{Consecución de los objetivos específicos del TFG.}
  \label{tab:objetivos_cumplidos}
  \begin{tabular}{|c|l|c|l|}
    \hline
    \textbf{Obj.} & \textbf{Descripción} & \textbf{Estado} & \textbf{Sección} \\
    \hline
    OE-1 & Rotulación, stream blanker y AFV  & Cumplido & \S\ref{sec:overlays}, \S\ref{sec:fuentes_auxiliares} \\
    OE-2 & Telemetría JSON + RabbitMQ        & Cumplido & \S\ref{sec:telemetria} \\
    OE-3 & Contenerización Docker Compose    & Cumplido & \S\ref{sec:docker} \\
    OE-4 & Despliegue en Kubernetes          & Cumplido & \S\ref{sec:kubernetes} \\
    OE-5 & Validación multi-escenario        & Cumplido & Cap.~\ref{chap:pruebas} \\
    OE-6 & Despliegue real en CyberNEMO      & Cumplido & \S\ref{sec:casos_uso_remi} \\
    \hline
  \end{tabular}
\end{table}

\vspace{0.3cm}

Más allá de la consecución individual de los objetivos, el trabajo ha demostrado que
es viable construir una plataforma de producción audiovisual profesional íntegramente
sobre software libre. La elección de GStreamer como motor multimedia ha resultado
especialmente acertada: su modelo de \textit{pipeline} modular ha permitido integrar
cada nueva funcionalidad como un módulo independiente sin necesidad de alterar el
núcleo original del sistema, lo que facilita el mantenimiento futuro y la contribución
por parte de terceros.

La separación entre \texttt{voctocore} y \texttt{voctogui} ha demostrado ser robusta
en todos los entornos probados y escala de forma natural desde una red local hasta
una infraestructura contenerizada y orquestada.
